\documentclass[11pt]{amsart}
\usepackage{geometry}                % See geometry.pdf to learn the layout options. There are lots.
\geometry{letterpaper}                   % ... or a4paper or a5paper or ... 
%\geometry{landscape}                % Activate for for rotated page geometry
%\usepackage[parfill]{parskip}    % Activate to begin paragraphs with an empty line rather than an indent
\usepackage{graphicx}
\usepackage{amssymb}
\usepackage{epstopdf}
\DeclareGraphicsRule{.tif}{png}{.png}{`convert #1 `dirname #1`/`basename #1 .tif`.png}

\title{Problema 3.14}

%\date{}                                           % Activate to display a given date or no date

\begin{document}
\maketitle
%\section{}
%\subsection{}
\noindent En un estadio se tienen 5 tipos diferentes de localidades, las cuales se identifican
por una clave numérica que es un valor comprendido entre 1 y 5. Los precios de
cada localidad y los datos referentes a las ventas de boletos para el próximo juego se proporcionan como sigue:
\newline

\noindent Datos: Pl, P2, P3, P4, P5

\noindent CLAVE1, CANT1

\noindent CLAVE2, CANT2

\noindent ...    ...

\noindent - 1 , - 1
\newline

\noindent Donde:
\newline

\noindent Pl, P2, P3,P4 y P5 son variables de tipo real que representan los precios de las localidades 1, 2, 3, 4 y 5, respectivamente.
\newline

\noindent CLAVE es una variable de tipo entero que representa el tipo de localidad de la venta i.
\newline

\noindent CANTi es una variable de tipo entero que indica la cantidad de boletos vendidos de un cierto tipo, en la venta i.
\newline

\noindent Construya un diagrama de flujo que:

a) Lea los precios.

b) Lea los datos de las ventas de boletos.

c) Imprima para cada venta, la clave, la cantidad y el importe total de los boletos 

vendidos en esta venta.

d) Calcule e imprima la cantidad de boletos vendidos de cada tipo.

e) Calcule e imprima la recaudación total del estadio.
\newline

\noindent\textbf{Explicación de las variables}
\newline

\noindent API, AP2, AP3, AP4 y AP5 : Variables de tipo entero. Acumulan el total de boletos
vendidos del tipo 1, 2, 3, 4 y 5, respectivamente.
\newline
\noindent RECAU: Variable de tipo real. Acumula la recaudación total delestadio.
\newline
\noindent P l, P2, P3, P4 y P5 : Variables de tipo real.
\newline
\noindent CLAVE y CANT: Variable de tipo entero.
\newline
\noindent PRE: Variable de tipo real. Almacena el total vendido en cada venta.
\newline
\noindent  A continuación en la siguiente tabla, podemos observar el seguimiento del
diagrama de flujo.
\newline

\noindent \begin{tabular}{|c|c|c|c|c|c|c|c|c|c|c|c|c|c|}
    \hline
    \multicolumn{14}{|c|}{\textbf{Tabla}} \\
    \hline
    \textbf{Pl} & \textbf{P2} & \textbf{P3} & \textbf{P4} & \textbf{P5} & \textbf{CLAVE} & \textbf{CANT} & \textbf{PRE} & \textbf{API} & \textbf{AP2} & \textbf{AP3} & \textbf{AP4} & \textbf{AP5} & \textbf{RECAU} \\
    \hline
    7.25 & 15.80 & 25.00 & 50.00 & 75.00 & 2 & 3 &  & 0 & 0 & 0 & 0 & 0 \\
    \hline
     & & & & & 3 & 8 & 47.40 & &3 &  & & &47.40    \\
    \hline
    & & & & &2 & 4 & 200.00 &  & &8 &  & & 247.40 \\
    \hline
    & & & & & 1 & 6 & 63.20 & & 7 & & & & 310.60  \\
    \hline
    & & & & & 4 & 5 & 43.50 & 6 & & & && 354.10  \\
    \hline
   & & & & & 1 & 12 & 250.00 & & && 5 & &604.10    \\
    \hline
    & & & & & 2 & 8 & 87.00 & 18 & & & && 691.10   \\
    \hline
    & & & & & 5 & 3 & 126.40 & &15 & & & & 817.50   \\
    \hline
    & & & & & 2 & 7 & 225.00 & & & & &3 & 1042.50  \\
    \hline
    & & & & & 3 & 14 & 110.60 &  &22 & & & &1153.10   \\
    \hline
    & & & & &  4 & 1 & 350.00 & & &22 & & &1503.10  \\
    \hline
    & & & & &2 & 11 & 50.00 & & & &6 && 1553.10   \\
    \hline
    & & & & &4 & 9 & 173.80 & &33 & & & &1726.90    \\
    \hline
    & & & & &5 & 7 & 450.00 & & & &15 & &  2176.90  \\
    \hline
    & & & & &1 & 23 & 525.00 & & & & &10 & 2701.90   \\
    \hline
   & & & & &  1 & 18 & 166.75 && 41 &  & & & 2868.65 \\
    \hline
    & & & & &3 & 4& 130.50 & 59& & & &10 & 2999.15   \\
    \hline
   & & & & &  -1 & -1 & 100.00 & &  & 26 & & & 3099.15 \\
    \hline

    
    \end{tabular}
\end{document}  